% ==========================================
% BAB VII PENUTUP
% ==========================================
\cleardoublepage
\chapter{PENUTUP}
\label{chap:penutup}

\section{Kesimpulan}
\label{sec:kesimpulan}

Tugas akhir ini telah merancang dan mengimplementasikan arsitektur sistem RAG-LLM untuk mitigasi halusinasi pada sistem tanya jawab biomedis dengan mengintegrasikan tiga teknik utama, yaitu ekspansi akronim, \textit{query rewriting} (QR), dan \textit{context reranking} (CR). Arsitektur tersebut disusun secara berlapis sehingga setiap teknik bekerja pada tahap yang berbeda di sepanjang alur RAG. Ekspansi akronim diterapkan pada tahap praproses korpus, sehingga seluruh dokumen yang masuk ke indeks sudah memuat kepanjangan dari akronim biomedis yang muncul. \textit{Hybrid retrieval} menggabungkan BM25 dan \textit{dense embedding} melalui \textit{Reciprocal Rank Fusion} sebagai mesin pencari ganda yang menutupi kelemahan satu sama lain. \textit{Query rewriting} memperkaya kueri pengguna di sisi masukan, dan \textit{context reranking} menyaring ulang urutan dokumen menggunakan \textit{cross-encoder} sebelum jawaban dihasilkan oleh LLM. Rancangan berlapis ini memungkinkan kontribusi tiap teknik diukur secara terpisah maupun gabungan, sekaligus memberi titik intervensi pada beberapa tahap yang berbeda di mana halusinasi dapat muncul.

Pada sisi implementasi, sistem yang dirancang telah diwujudkan secara utuh dan bekerja \textit{end-to-end} pada 500 pertanyaan PubMedQA dengan empat pilihan LLM, yaitu GPT-4.1-mini, Claude Haiku 4.5, Llama 3.3 70B, dan Llama 3.2 3B lokal. Korpus dipersiapkan menggunakan algoritma Schwartz-Hearst untuk mengekstrak pasangan akronim dan kepanjangannya, lalu indeks BM25 serta basis vektor Chroma dibangun di atas korpus yang sudah diperkaya. Pada tahap \textit{runtime}, pertanyaan pengguna ditulis ulang oleh LLM, ditarik dokumennya oleh \textit{hybrid retriever}, diurutkan ulang oleh \textit{cross-encoder}, lalu lima dokumen teratas dirangkai bersama \textit{prompt} dan dikirim ke LLM untuk menghasilkan jawaban. Seluruh komponen tersebut berhasil terintegrasi pada empat pilihan LLM tanpa modifikasi besar pada masing-masing model, yang menunjukkan bahwa arsitektur mitigasi yang dirancang bersifat \textit{model-agnostic} dan dapat diterapkan pada berbagai \textit{backbone} LLM. Dengan demikian, rumusan masalah dan tujuan pertama mengenai perancangan dan implementasi sistem telah terjawab.

Menjawab rumusan masalah kedua mengenai kontribusi komparatif setiap teknik mitigasi, analisis terhadap dua belas konfigurasi memberi tiga temuan utama mengenai kontribusi setiap teknik terhadap mitigasi halusinasi. \textbf{Pertama}, ekspansi akronim terbukti memberi dampak positif yang paling konsisten dalam menurunkan halusinasi. Penerapannya menaikkan \textit{accuracy} sebesar +2{,}2 hingga +4{,}8 poin persentase pada keempat konfigurasi mitigasi, dan yang lebih mencolok adalah kenaikan \textit{faithfulness} rata-rata sebesar +9 poin persentase. Peningkatan \textit{faithfulness} menandakan bahwa jawaban yang dihasilkan LLM menjadi lebih setia pada konteks dokumen, yang merupakan indikator langsung berkurangnya halusinasi intrinsik. \textbf{Kedua}, \textit{query rewriting} yang berdiri sendiri ternyata tidak memberi peningkatan yang berarti pada \textit{accuracy} maupun pada metrik RAGAS, dan bahkan pada beberapa konfigurasi sedikit menurunkan \textit{recall} karena tulis-ulang kueri menggeser kata kunci yang sebelumnya berhasil ditangkap oleh BM25. Namun ketika \textit{query rewriting} digabungkan dengan \textit{context reranking}, kombinasi keduanya justru menghasilkan konfigurasi terbaik, yaitu \textit{Hybrid Retrieval} dengan ekspansi akronim, QR, dan CR, yang mencapai \textit{accuracy} 74{,}8\% dan \textit{faithfulness} 0{,}980. Konfigurasi ini juga menghasilkan metrik \textit{retrieval} paling tinggi pada seluruh dua belas konfigurasi (Recall@5 87{,}08\%, MAP@5 84{,}53\%, dan nDCG@5 90{,}11\%). \textbf{Ketiga}, \textit{context reranking} terbukti konsisten menurunkan halusinasi pada ketiga kelompok konfigurasi. \textit{Cross-encoder} yang membaca pasangan (kueri, dokumen) secara bersamaan dapat memperbaiki urutan dokumen yang awalnya kurang tepat dari \textit{retriever} cepat, sehingga LLM menerima konteks yang lebih relevan dan kecenderungan menghasilkan jawaban yang tidak \textit{grounded} berkurang. Dampak positif ini terlihat baik pada metrik \textit{retrieval} maupun pada metrik \textit{generation}.

Secara keseluruhan, konfigurasi terbaik yang menggabungkan keempat komponen, yaitu \textit{Hybrid Retrieval}, ekspansi akronim, \textit{query rewriting}, dan \textit{context reranking}, berhasil meningkatkan \textit{accuracy} dari 65{,}4\% pada \textit{baseline} BM25 paling sederhana menjadi 74{,}8\%, dengan \textit{faithfulness} mencapai 0{,}980. Sejalan dengan itu, sistem juga berhasil menurunkan tingkat halusinasi sebesar 9{,}4 poin persentase, yaitu dari 34{,}6\% pada \textit{baseline} BM25 menjadi 25{,}2\% pada konfigurasi terbaik. Hasil ini menegaskan bahwa pendekatan mitigasi berlapis pada tiga titik berbeda di alur RAG, yaitu di praproses korpus, di sisi kueri, dan di tahap pemeringkatan akhir, mampu menurunkan halusinasi secara substansial pada sistem tanya jawab biomedis dan dapat dijadikan dasar pengembangan sistem produksi pada domain serupa.

\section{Saran}
\label{sec:saran}

Sistem yang dibangun pada tugas akhir ini sudah menunjukkan hasil yang cukup baik, tetapi masih banyak hal yang dapat dikerjakan lebih lanjut agar mitigasi halusinasi pada RAG-LLM biomedis semakin matang. Berikut beberapa arah pengembangan yang dapat dipertimbangkan.

\begin{enumerate}
    \item \textbf{Eksplorasi \textit{prompt engineering} yang lebih mendalam.}
    \textit{Prompt} jawaban dan \textit{prompt} \textit{query rewriting} yang dipakai pada tugas akhir ini didapat melalui beberapa kali percobaan sederhana hingga didapat versi yang stabil. Namun masih banyak teknik \textit{prompt} dalam literatur yang belum sempat diuji, misalnya \textit{chain-of-thought}, \textit{few-shot} dengan contoh biomedis, atau \textit{self-consistency decoding}. Pengujian teknik-teknik tersebut secara lebih sistematis layak dilakukan karena peluang untuk menaikkan kualitas jawaban dan \textit{faithfulness} masih terbuka tanpa harus mengubah arsitektur \textit{retrieval} yang sudah ada.

    \item \textbf{Evaluasi langsung oleh pakar biomedis.}
    Metrik seperti \textit{accuracy} dan \textit{faithfulness} memang berguna untuk membandingkan banyak konfigurasi sekaligus, tetapi angka saja belum cukup menggambarkan apakah jawaban yang dihasilkan benar-benar aman bila dibaca oleh tenaga medis maupun pasien. Penilaian langsung dari pakar biomedis dapat menangkap aspek-aspek yang sulit diukur oleh metrik otomatis, seperti kebenaran klinis, kemungkinan salah tafsir, atau informasi penting yang luput dijawab. Evaluasi semacam ini menjadi semakin penting bila sistem diarahkan ke konteks penggunaan nyata, baik untuk membantu edukasi pasien maupun mendukung pekerjaan tenaga medis sehari-hari.

    \item \textbf{Perluasan cakupan korpus dari \textit{paper} ilmiah biomedis.}
    Korpus pada tugas akhir ini terbatas pada subset \textit{pqa\_labeled} dari PubMedQA yang berisi seribu \textit{paper}. Untuk menjangkau lebih banyak topik medis, sumber data dapat diperluas ke subset PubMed yang lebih besar, basis data lain seperti BioASQ atau MedQA, atau bahkan catatan medis seperti MIMIC. Penambahan ini akan memperkaya variasi pertanyaan yang dapat ditangani sistem dan membuat hasil evaluasi lebih dekat ke kebutuhan klinis sebenarnya.

    \item \textbf{Eksperimen kombinasi metode optimalisasi RAG-LLM lainnya.}
    Literatur mitigasi halusinasi pada RAG-LLM berkembang dengan sangat cepat. Survei pada jurnal seperti \textit{Mathematics} serta \textit{paper} terbaru lain menyebut banyak teknik yang menjanjikan tetapi belum sempat dimasukkan ke dalam tugas akhir ini, seperti \textit{iterative retrieval}, \textit{Hypothetical Document Embeddings} (HyDE), \textit{self-RAG}, \textit{Corrective RAG}, dan \textit{constrained decoding}. Eksperimen yang membandingkan atau menggabungkan teknik-teknik tersebut dengan arsitektur yang sudah dibangun akan memberi gambaran tentang batas atas kinerja yang masih dapat diraih pada domain biomedis.

    \item \textbf{Pengujian pada model LLM generasi yang lebih baru.}
    Model bahasa berkembang dengan kecepatan tinggi, sehingga menarik bila pendekatan ini diuji ulang pada generasi model berikutnya seperti GPT-5, Claude Opus 5, atau Llama 4. Hasil pengujian tersebut dapat menunjukkan apakah lapisan mitigasi eksplisit yang dirancang pada tugas akhir ini masih dibutuhkan, atau apakah kemampuan \textit{grounding} bawaan pada model yang lebih kuat sudah cukup untuk menggantikan sebagian perannya. Pengujian pada model \textit{open-source} berukuran menengah yang lebih murah secara komputasi juga relevan, agar arsitektur ini tetap dapat dijalankan pada lingkungan dengan anggaran terbatas.

    \item \textbf{Dukungan multibahasa dan adaptasi domain.}
    Tugas akhir ini dibatasi pada pertanyaan biomedis berbahasa Inggris. Pengembangan ke arah \textit{multilingual biomedical question answering} akan memperluas jangkauan pengguna, sementara adaptasi domain (\textit{domain adaptation}) ke subbidang medis yang lebih spesifik, misalnya onkologi atau kardiologi, dapat meningkatkan ketepatan jawaban pada konteks yang lebih sempit.

    \item \textbf{Pelatihan khusus \textit{retriever} dan deteksi halusinasi aktif.}
    Kinerja \textit{retrieval} masih dapat ditingkatkan dengan melatih khusus (\textit{fine-tuning}) model \textit{retriever} pada korpus biomedis sehingga representasi istilah medis menjadi lebih tepat. Selain itu, penambahan deteksi halusinasi aktif (\textit{active hallucination detection}) yang memeriksa kesetiaan jawaban terhadap konteks secara langsung saat \textit{generation} dapat menjadi lapisan pengaman tambahan sebelum jawaban diberikan kepada pengguna.
\end{enumerate}

Beberapa arah pengembangan tersebut dapat menjadi langkah lanjutan dari tugas akhir ini agar sistem tanya jawab biomedis berbasis RAG-LLM semakin dekat dengan kebutuhan klinis nyata.
